%%%%%%%%%%%%%%%%%%%%%%%%%%%%%%%%%%%%%%%%%%%%%%%%%%%%%%%%%%%%%%%%%%%
% Physik IV - Blatt 06 Lernfassung / Abgabe / Tafelvortrag
% Standalone university LaTeX template -- one-file version
%%%%%%%%%%%%%%%%%%%%%%%%%%%%%%%%%%%%%%%%%%%%%%%%%%%%%%%%%%%%%%%%%%%

\documentclass[11pt,a4paper]{scrartcl}

%%%%%%%%%%%%%%%%%%%%%%%%%%%%%%%%%%%%%%%%%%%%%%%%%%%%%%%%%%%%%%%%%%%
% Embedded file: includes/university_metadata.tex
%%%%%%%%%%%%%%%%%%%%%%%%%%%%%%%%%%%%%%%%%%%%%%%%%%%%%%%%%%%%%%%%%%%
\newcommand{\CourseTitle}{COURSE NAME}
\newcommand{\CourseShortTitle}{COURSE SHORT NAME}
\newcommand{\DocumentKind}{DOCUMENT KIND}
\newcommand{\DocumentTitle}{DOCUMENT TITLE}
\newcommand{\DocumentSubtitle}{}
\newcommand{\AuthorName}{YOUR NAME}
\newcommand{\InstitutionName}{INSTITUTION NAME}
\newcommand{\SubmissionDate}{\today}

\newcommand{\makeuniversitytitle}{%
    \begin{titlepage}
        \centering
        \vspace*{2cm}
        {\Large \CourseTitle\par}
        \vspace{1.2cm}
        {\huge\bfseries \DocumentTitle\par}
        \ifx\DocumentSubtitle\empty\else
            \vspace{0.4cm}
            {\Large \DocumentSubtitle\par}
        \fi
        \vspace{1cm}
        {\Large \DocumentKind\par}
        \vfill
        {\large \AuthorName\par}
        \vspace{0.2cm}
        {\large \InstitutionName\par}
        \vspace{0.8cm}
        {\large \SubmissionDate\par}
    \end{titlepage}
    \pagenumbering{arabic}
}

\newcommand{\makechaptertitle}{%
    \begin{center}
        {\Large \CourseTitle\par}
        \vspace{0.5cm}
        {\huge\bfseries \DocumentTitle\par}
        \ifx\DocumentSubtitle\empty\else
            \vspace{0.25cm}
            {\Large \DocumentSubtitle\par}
        \fi
        \vspace{0.5cm}
        {\large \AuthorName \hfill \SubmissionDate\par}
    \end{center}
    \vspace{0.5cm}
}

\newcommand{\makephysiktitle}{\makeuniversitytitle}

%%%%%%%%%%%%%%%%%%%%%%%%%%%%%%%%%%%%%%%%%%%%%%%%%%%%%%%%%%%%%%%%%%%
% Document metadata
%%%%%%%%%%%%%%%%%%%%%%%%%%%%%%%%%%%%%%%%%%%%%%%%%%%%%%%%%%%%%%%%%%%
\renewcommand{\CourseTitle}{Physik IV: Kerne und Teilchen}
\renewcommand{\CourseShortTitle}{Physik IV}
\renewcommand{\DocumentKind}{Hausaufgabe, Tafelvortrag und Lernfassung}
\renewcommand{\DocumentTitle}{Aufgabenblatt 6}
\renewcommand{\DocumentSubtitle}{Spin-Bahn-Wechselwirkung und Schalenmodell}
\renewcommand{\AuthorName}{Tobias Laser}
\renewcommand{\InstitutionName}{Universit\"at Innsbruck}
\renewcommand{\SubmissionDate}{28. April 2026}

%%%%%%%%%%%%%%%%%%%%%%%%%%%%%%%%%%%%%%%%%%%%%%%%%%%%%%%%%%%%%%%%%%%
% Embedded file: includes/university_packages.tex
%%%%%%%%%%%%%%%%%%%%%%%%%%%%%%%%%%%%%%%%%%%%%%%%%%%%%%%%%%%%%%%%%%%
\usepackage[margin=2.5cm]{geometry}
\usepackage[onehalfspacing]{setspace}
\usepackage{scrlayer-scrpage}

\usepackage[T1]{fontenc}
\usepackage[utf8]{inputenc}
\usepackage[ngerman]{babel}
\usepackage{lmodern}
\usepackage{microtype}
\usepackage{csquotes}

\usepackage{amsmath,amsthm,amssymb,mathtools}
\usepackage{bm}
\usepackage{icomma}
\usepackage{cancel}

\usepackage{graphicx}
\usepackage{tikz}
\usetikzlibrary{arrows.meta,calc,decorations.pathmorphing,positioning,patterns,angles,quotes}
\usepackage{pgfplots}
\pgfplotsset{compat=1.18}

\usepackage[locale=DE,space-before-unit=true,per-mode=symbol,separate-uncertainty=true]{siunitx}
\usepackage{booktabs,multirow,array,tabularx,longtable}
\usepackage{caption,subcaption}
\usepackage{placeins}
\usepackage{enumitem}
\usepackage{xparse}

\usepackage[most]{tcolorbox}
\tcbuselibrary{breakable,skins,theorems}

\usepackage[breaklinks=true,colorlinks=true,linkcolor=blue,urlcolor=blue,citecolor=blue]{hyperref}
\usepackage[nameinlink,noabbrev]{cleveref}

\clearpairofpagestyles
\ihead{\CourseShortTitle}
\ohead{\DocumentKind}
\cfoot{\pagemark}
\pagestyle{scrheadings}

\setlist[enumerate]{itemsep=0.4em,topsep=0.4em}
\setlist[itemize]{itemsep=0.25em,topsep=0.35em}
\captionsetup{font=small,labelfont=bf}

%%%%%%%%%%%%%%%%%%%%%%%%%%%%%%%%%%%%%%%%%%%%%%%%%%%%%%%%%%%%%%%%%%%
% Embedded file: includes/university_macros.tex
%%%%%%%%%%%%%%%%%%%%%%%%%%%%%%%%%%%%%%%%%%%%%%%%%%%%%%%%%%%%%%%%%%%
\newcommand{\R}{\mathbb{R}}
\newcommand{\N}{\mathbb{N}}
\newcommand{\Z}{\mathbb{Z}}
\newcommand{\Q}{\mathbb{Q}}
\newcommand{\C}{\mathbb{C}}

\newcommand{\set}[1]{\left\{ #1 \right\}}
\newcommand{\abs}[1]{\left| #1 \right|}
\newcommand{\norm}[1]{\left\lVert #1 \right\rVert}
\newcommand{\avg}[1]{\left\langle #1 \right\rangle}
\newcommand{\paren}[1]{\left( #1 \right)}
\newcommand{\bracket}[1]{\left[ #1 \right]}
\newcommand{\ol}[1]{\overline{#1}}

\newcommand{\dd}{\,\mathrm{d}}
\newcommand{\dv}[2]{\frac{\dd #1}{\dd #2}}
\newcommand{\pdv}[2]{\frac{\partial #1}{\partial #2}}
\newcommand{\pddv}[2]{\frac{\partial^2 #1}{\partial #2^2}}
\newcommand{\evalat}[2]{\left. #1 \right|_{#2}}

\newcommand{\vect}[1]{\bm{#1}}
\newcommand{\unitvect}[1]{\hat{\bm{#1}}}
\newcommand{\grad}{\vect{\nabla}}
\newcommand{\divergence}{\grad\!\cdot}
\newcommand{\curl}{\grad\!\times}

\newcommand{\half}{\frac{1}{2}}
\newcommand{\third}{\frac{1}{3}}
\newcommand{\twothirds}{\frac{2}{3}}
\newcommand{\hbarc}{\hbar c}
\newcommand{\alphaf}{\alpha}
\newcommand{\eV}{\si{eV}}
\newcommand{\MeV}{\si{MeV}}
\newcommand{\GeV}{\si{GeV}}
\newcommand{\TeV}{\si{TeV}}

\newcommand{\nuclide}[3]{^{#1}_{#2}\mathrm{#3}}
\newcommand{\isotope}[2]{^{#1}\mathrm{#2}}
\newcommand{\anti}[1]{\overline{#1}}
\newcommand{\pbar}{\anti{p}}
\newcommand{\nbar}{\anti{n}}
\newcommand{\ee}{e^-}
\newcommand{\ep}{e^+}
\newcommand{\mup}{\mu^+}
\newcommand{\mum}{\mu^-}
\newcommand{\nue}{\nu_e}
\newcommand{\numu}{\nu_\mu}
\newcommand{\nutau}{\nu_\tau}
\newcommand{\pionp}{\pi^+}
\newcommand{\pionm}{\pi^-}
\newcommand{\pionz}{\pi^0}
\newcommand{\kaonp}{K^+}
\newcommand{\kaonm}{K^-}
\newcommand{\kaonz}{K^0}

\newcommand{\diffxs}{\frac{\dd\sigma}{\dd\Omega}}
\newcommand{\totalxs}{\sigma_{\mathrm{total}}}
\newcommand{\lum}{\mathcal{L}}
\newcommand{\smand}{s}
\newcommand{\tmand}{t}
\newcommand{\umand}{u}
\newcommand{\Ecm}{E_{\mathrm{CM}}}

\newcommand{\ie}{d.\,h.}
\newcommand{\eg}{z.\,B.}
\newcommand{\wrt}{bez\"uglich}

\DeclareSIUnit{\barn}{b}
\DeclareSIUnit{\millibarn}{mb}
\DeclareSIUnit{\microbarn}{\micro b}
\DeclareSIUnit{\nanobarn}{nb}
\DeclareSIUnit{\fermi}{fm}
\DeclareSIUnit{\year}{a}

%%%%%%%%%%%%%%%%%%%%%%%%%%%%%%%%%%%%%%%%%%%%%%%%%%%%%%%%%%%%%%%%%%%
% Embedded file: includes/university_boxes.tex
%%%%%%%%%%%%%%%%%%%%%%%%%%%%%%%%%%%%%%%%%%%%%%%%%%%%%%%%%%%%%%%%%%%
\tcbset{
    unibox/.style={
        enhanced,
        breakable,
        sharp corners,
        boxrule=0.6pt,
        left=1.2mm,
        right=1.2mm,
        top=1mm,
        bottom=1mm,
        before skip=0.8em,
        after skip=0.8em,
        fonttitle=\bfseries
    }
}

\newtcolorbox{calculationbox}[1][]{unibox,title=Rechnung,colback=black!2,colframe=black!55,#1}
\newtcolorbox{sidecalcbox}[1][]{unibox,title=Nebenrechnung,colback=black!1,colframe=black!40,#1}
\newtcolorbox{solutionbox}[1][]{unibox,title=L\"osung,colback=blue!2,colframe=blue!45!black,#1}
\newtcolorbox{answerbox}[1][]{unibox,title=Antwort,colback=green!3,colframe=green!45!black,#1}
\newtcolorbox{notebox}[1][]{unibox,title=Notiz,colback=yellow!6,colframe=yellow!45!black,#1}
\newtcolorbox{definitionbox}[1][]{unibox,title=Definition,colback=cyan!3,colframe=cyan!45!black,#1}
\newtcolorbox{warningbox}[1][]{unibox,title=Achtung,colback=red!3,colframe=red!55!black,#1}
\newtcolorbox{summarybox}[1][]{unibox,title=Zusammenfassung,colback=violet!3,colframe=violet!50!black,#1}
\newtcolorbox{formulabox}[1][]{unibox,title=Formel,colback=orange!4,colframe=orange!65!black,#1}
\newtcolorbox{taskbox}[1][]{unibox,title=Aufgabe,colback=white,colframe=black!75,#1}
\newtcolorbox{talkbox}[1][]{unibox,title=Tafelvortrag,colback=teal!3,colframe=teal!50!black,#1}
\newtcolorbox{learningbox}[1][]{unibox,title=Lernidee,colback=purple!3,colframe=purple!50!black,#1}

%%%%%%%%%%%%%%%%%%%%%%%%%%%%%%%%%%%%%%%%%%%%%%%%%%%%%%%%%%%%%%%%%%%
% Embedded file: includes/university_theorems.tex
%%%%%%%%%%%%%%%%%%%%%%%%%%%%%%%%%%%%%%%%%%%%%%%%%%%%%%%%%%%%%%%%%%%
\theoremstyle{definition}
\newtheorem{definition}{Definition}[section]
\newtheorem{example}{Beispiel}[section]

\theoremstyle{plain}
\newtheorem{satz}{Satz}[section]
\newtheorem{lemma}{Lemma}[section]

\theoremstyle{remark}
\newtheorem*{remark}{Bemerkung}

%%%%%%%%%%%%%%%%%%%%%%%%%%%%%%%%%%%%%%%%%%%%%%%%%%%%%%%%%%%%%%%%%%%
\begin{document}

\makeuniversitytitle
\tableofcontents
\newpage

%%%%%%%%%%%%%%%%%%%%%%%%%%%%%%%%%%%%%%%%%%%%%%%%%%%%%%%%%%%%%%%%%%%
\section*{Einordnung des Blattes}
\addcontentsline{toc}{section}{Einordnung des Blattes}

\begin{summarybox}[title=Ziel dieses Blattes]
Dieses Blatt verbindet zwei Kernideen des nuklearen Schalenmodells:
\begin{itemize}
    \item Die \textbf{Spin-Bahn-Wechselwirkung} spaltet ein Niveau mit Bahndrehimpuls $l$ in zwei Niveaus mit $j=l\pm \half$ auf.
    \item Aus \textbf{Bindungsenergien} benachbarter Kerne kann man Einteilchenenergien bzw. Schalenabst\"ande absch\"atzen.
\end{itemize}
Der wichtigste konzeptionelle Punkt ist: In der Kernphysik liegt das Niveau mit $j=l+\half$ experimentell tiefer als das Niveau mit $j=l-\half$. Das ist f\"ur die Entstehung der magischen Zahlen entscheidend.
\end{summarybox}

\begin{notebox}[title=Benutzte Konstanten und Konventionen]
\[
    \alpha \hbar c \approx \SI{1.440}{MeV.fm},
    \qquad
    R_0 \approx \SI{1.2}{fm}.
\]
Wenn nicht anders angegeben, werden Drehimpulse in Einheiten von $\hbar$ geschrieben. Dann gilt z.\,B. $\avg{\vect{j}^{\,2}}=j(j+1)$ statt $\hbar^2 j(j+1)$.
\end{notebox}

\FloatBarrier
\newpage

%%%%%%%%%%%%%%%%%%%%%%%%%%%%%%%%%%%%%%%%%%%%%%%%%%%%%%%%%%%%%%%%%%%
\section{Aufgabe 1: Spin-Bahn-Wechselwirkung}

\begin{taskbox}
Gegeben ist der Korrekturterm der Spin-Bahn-Wechselwirkung
\[
    V_{ls}(r,\vect{l},\vect{s}) = V_{ls}(r)\,\paren{\vect{l}\cdot\vect{s}}.
\]
Die explizite radiale Form von $V_{ls}(r)$ soll nicht verwendet werden. Entscheidend ist nur der Operator $\vect{l}\cdot\vect{s}$.
\begin{enumerate}[label=(\alph*)]
    \item Zeigen Sie mit Hilfe von $\vect{j}=\vect{l}+\vect{s}$:
    \[
        \vect{l}\cdot\vect{s}
        = \half \paren{\vect{j}^{\,2}-\vect{l}^{\,2}-\vect{s}^{\,2}}.
    \]
    \item Betrachten Sie die $1d$-Schale. Welchen Bahndrehimpuls $l$ besitzt diese? In welche Werte von $j$ wird sie aufgespalten?
    \item Berechnen Sie $\avg{\vect{l}\cdot\vect{s}}$ f\"ur beide Unterschalen der $1d$-Schale und geben Sie die Energiedifferenz in Abh\"angigkeit von $V_{ls}(r)$ an.
\end{enumerate}
\end{taskbox}

\subsection{Lernteil: Was bedeutet Spin-Bahn-Wechselwirkung?}

\begin{learningbox}
Im Schalenmodell bewegen sich Protonen und Neutronen in einem mittleren Kernpotential. Ein Zustand wird nicht nur durch den Bahndrehimpuls $l$, sondern durch Kopplung von Bahndrehimpuls und Spin durch den Gesamtdrehimpuls
\[
    \vect{j}=\vect{l}+\vect{s}
\]
beschrieben. Da Nukleonen Spin $s=\half$ besitzen, entstehen f\"ur jedes $l>0$ zwei M\"oglichkeiten:
\[
    j=l+\half,
    \qquad
    j=l-\half.
\]
Die Spin-Bahn-Wechselwirkung macht diese beiden M\"oglichkeiten energetisch verschieden. Dadurch wird aus einer $nl$-Schale eine Aufspaltung in $nl_{j=l+1/2}$ und $nl_{j=l-1/2}$.
\end{learningbox}

\subsection{Abgabe: Ausf\"uhrliche L\"osung}

\begin{enumerate}[label=(\alph*)]
    \item \textbf{Herleitung des Operators $\vect{l}\cdot\vect{s}$.}

    \begin{calculationbox}
        Ausgangspunkt ist die Definition des Gesamtdrehimpulses:
        \[
            \vect{j}=\vect{l}+\vect{s}.
        \]
        Quadrieren liefert
        \[
            \vect{j}^{\,2}
            =\paren{\vect{l}+\vect{s}}^2.
        \]
        Das Skalarprodukt wird ausmultipliziert:
        \[
            \paren{\vect{l}+\vect{s}}^2
            =\vect{l}^{\,2}+2\vect{l}\cdot\vect{s}+\vect{s}^{\,2}.
        \]
        Also gilt
        \[
            \vect{j}^{\,2}
            =\vect{l}^{\,2}+2\vect{l}\cdot\vect{s}+\vect{s}^{\,2}.
        \]
        Umstellen nach $\vect{l}\cdot\vect{s}$ ergibt
        \[
            2\vect{l}\cdot\vect{s}
            =\vect{j}^{\,2}-\vect{l}^{\,2}-\vect{s}^{\,2}
        \]
        und damit
        \[
            \boxed{\vect{l}\cdot\vect{s}
            =\half\paren{\vect{j}^{\,2}-\vect{l}^{\,2}-\vect{s}^{\,2}}}.
        \]
    \end{calculationbox}

    \begin{answerbox}
        Die gesuchte Identit\"at folgt direkt aus dem Quadrieren von $\vect{j}=\vect{l}+\vect{s}$.
    \end{answerbox}

    \item \textbf{Aufspaltung der $1d$-Schale.}

    \begin{calculationbox}
        In der spektroskopischen Notation steht der Buchstabe f\"ur den Bahndrehimpuls:
        \[
            s\to l=0,
            \qquad
            p\to l=1,
            \qquad
            d\to l=2,
            \qquad
            f\to l=3.
        \]
        F\"ur die $1d$-Schale gilt daher
        \[
            l=2.
        \]
        Der Spin eines Nukleons ist
        \[
            s=\half.
        \]
        Bei der Kopplung von $l$ und $s$ sind allgemein die Werte
        \[
            j=\abs{l-s},\abs{l-s}+1,\dots,l+s
        \]
        erlaubt. Da $s=\half$ ist, bleiben nur
        \[
            j=l-\half
            \qquad\text{und}\qquad
            j=l+\half.
        \]
        Mit $l=2$ folgt
        \[
            j=2-\half=\frac{3}{2},
            \qquad
            j=2+\half=\frac{5}{2}.
        \]
    \end{calculationbox}

    \begin{answerbox}
        Die $1d$-Schale besitzt $l=2$ und spaltet in
        \[
            \boxed{1d_{5/2}\quad\text{und}\quad 1d_{3/2}}
        \]
        auf.
    \end{answerbox}

    \item \textbf{Erwartungswerte und Energiedifferenz.}

    \begin{calculationbox}
        Aus Teil (a) folgt f\"ur Eigenzust\"ande von $j^2,l^2,s^2$:
        \[
            \avg{\vect{l}\cdot\vect{s}}
            =\half\bracket{j(j+1)-l(l+1)-s(s+1)}.
        \]
        F\"ur die $1d$-Schale gilt
        \[
            l=2,
            \qquad
            s=\half.
        \]
        Damit
        \[
            l(l+1)=2\cdot 3=6,
            \qquad
            s(s+1)=\half\cdot\frac{3}{2}=\frac{3}{4}.
        \]
    \end{calculationbox}

    \begin{calculationbox}[title=Unterschale $1d_{5/2}$]
        F\"ur $j=\frac{5}{2}$ ist
        \[
            j(j+1)=\frac{5}{2}\cdot\frac{7}{2}=\frac{35}{4}.
        \]
        Daher
        \[
            \avg{\vect{l}\cdot\vect{s}}_{5/2}
            =\half\bracket{\frac{35}{4}-6-\frac{3}{4}}
            =\half\bracket{\frac{35-24-3}{4}}
            =\half\cdot\frac{8}{4}
            =1.
        \]
    \end{calculationbox}

    \begin{calculationbox}[title=Unterschale $1d_{3/2}$]
        F\"ur $j=\frac{3}{2}$ ist
        \[
            j(j+1)=\frac{3}{2}\cdot\frac{5}{2}=\frac{15}{4}.
        \]
        Daher
        \[
            \avg{\vect{l}\cdot\vect{s}}_{3/2}
            =\half\bracket{\frac{15}{4}-6-\frac{3}{4}}
            =\half\bracket{\frac{15-24-3}{4}}
            =\half\cdot\paren{-\frac{12}{4}}
            =-\frac{3}{2}.
        \]
    \end{calculationbox}

    \begin{calculationbox}[title=Energieverschiebungen]
        Die Energieverschiebung ist proportional zu
        \[
            \Delta E_j = V_{ls}(r)\,\avg{\vect{l}\cdot\vect{s}}_j.
        \]
        Somit
        \[
            \Delta E_{5/2}=V_{ls}(r)\cdot 1=V_{ls}(r),
        \]
        und
        \[
            \Delta E_{3/2}=V_{ls}(r)\cdot\paren{-\frac{3}{2}}
            =-\frac{3}{2}V_{ls}(r).
        \]
        Die Differenz in dieser Reihenfolge ist
        \[
            E_{5/2}-E_{3/2}
            =V_{ls}(r)-\paren{-\frac{3}{2}V_{ls}(r)}
            =\frac{5}{2}V_{ls}(r).
        \]
    \end{calculationbox}

    \begin{answerbox}
        \[
            \boxed{\avg{\vect{l}\cdot\vect{s}}_{1d_{5/2}}=1},
            \qquad
            \boxed{\avg{\vect{l}\cdot\vect{s}}_{1d_{3/2}}=-\frac{3}{2}}.
        \]
        Damit gilt
        \[
            \boxed{E_{5/2}-E_{3/2}=\frac{5}{2}V_{ls}(r)}.
        \]
        Experimentell ist in der Kernphysik $V_{ls}(r)$ effektiv so, dass $j=l+\half$ tiefer liegt als $j=l-\half$.
    \end{answerbox}
\end{enumerate}

\subsection{Tafelvortrag zu Aufgabe 1}

\begin{talkbox}
\textbf{Tafelstruktur in 4 Schritten:}
\begin{enumerate}
    \item Start mit $\vect{j}=\vect{l}+\vect{s}$ und durch Quadrieren direkt
    \[
        \vect{l}\cdot\vect{s}=\half\paren{\vect{j}^{\,2}-\vect{l}^{\,2}-\vect{s}^{\,2}}.
    \]
    \item F\"ur $1d$: Buchstabe $d$ bedeutet $l=2$; Nukleonenspin $s=\half$.
    \item Daher $j=l\pm\half=\frac{5}{2},\frac{3}{2}$, also $1d_{5/2}$ und $1d_{3/2}$.
    \item Erwartungswerte einsetzen:
    \[
        \avg{\vect{l}\cdot\vect{s}}
        =\half\bracket{j(j+1)-l(l+1)-s(s+1)}.
    \]
    Ergebnis:
    \[
        1d_{5/2}: 1,
        \qquad
        1d_{3/2}: -\frac{3}{2},
        \qquad
        E_{5/2}-E_{3/2}=\frac{5}{2}V_{ls}(r).
    \]
\end{enumerate}
\end{talkbox}

\FloatBarrier
\newpage

%%%%%%%%%%%%%%%%%%%%%%%%%%%%%%%%%%%%%%%%%%%%%%%%%%%%%%%%%%%%%%%%%%%
\section{Aufgabe 2: Schalenmodell}

\begin{taskbox}
Der Kern $\isotope{16}{O}$ ist doppeltmagisch, da er aus 8 Protonen und 8 Neutronen besteht. Gegeben sind die totalen Bindungsenergien
\[
    E_B(\isotope{15}{O})=\SI{111.96}{MeV},
    \qquad
    E_B(\isotope{16}{O})=\SI{127.62}{MeV},
    \qquad
    E_B(\isotope{17}{O})=\SI{131.76}{MeV},
\]
und
\[
    E_B(\isotope{17}{F})=\SI{128.22}{MeV}.
\]
\begin{enumerate}[label=(\alph*)]
    \item Berechnen Sie den Abstand zwischen den Neutronenschalen $1p_{1/2}$ und $1d_{5/2}$ f\"ur Kerne im Bereich $A\approx 16$.
    \item Interpretieren Sie den Unterschied in der totalen Bindungsenergie von $\isotope{17}{O}$ und $\isotope{17}{F}$.
    \item Sch\"atzen Sie aus $E_B(\isotope{17}{O})-E_B(\isotope{17}{F})$ den Radius ab. Verwenden Sie
    \[
        E_C=\frac{3}{5}\frac{Z(Z-1)\alpha\hbar c}{R}.
    \]
    \item Vergleichen Sie mit $R\approx R_0A^{1/3}$ f\"ur $R_0=\SI{1.2}{fm}$ und diskutieren Sie die Abweichung.
\end{enumerate}
\end{taskbox}

\subsection{Lernteil: Bindungsenergien als Schalenabst\"ande}

\begin{definitionbox}[title=Trennenergie]
Die Neutronentrennenergie eines Kerns ist n\"aherungsweise die Energie, die man aufbringen muss, um das zuletzt gebundene Neutron zu entfernen:
\[
    S_n(A,Z)=E_B(A,Z)-E_B(A-1,Z).
\]
F\"ur Protonen analog:
\[
    S_p(A,Z)=E_B(A,Z)-E_B(A-1,Z-1).
\]
Gro\ss{}e Trennenergie bedeutet: Das letzte Nukleon ist stark gebunden.
\end{definitionbox}

\begin{learningbox}
$\isotope{16}{O}$ besitzt $Z=N=8$. Im Schalenmodell ist damit die $1p_{1/2}$-Schale gerade voll. Entfernt man aus $\isotope{16}{O}$ ein Neutron, entfernt man ein Neutron aus $1p_{1/2}$. F\"ugt man zu $\isotope{16}{O}$ ein Neutron hinzu, kommt es in die n\"achste freie Schale $1d_{5/2}$. Deshalb kann man die L\"ucke zwischen $1p_{1/2}$ und $1d_{5/2}$ aus zwei Neutronentrennenergien absch\"atzen.
\end{learningbox}

\subsection{Abgabe: Ausf\"uhrliche L\"osung}

\begin{enumerate}[label=(\alph*)]
    \item \textbf{Abstand zwischen $1p_{1/2}$ und $1d_{5/2}$.}

    \begin{calculationbox}
        Zuerst entfernen wir gedanklich ein Neutron aus $\isotope{16}{O}$. Dann entsteht $\isotope{15}{O}$. Die daf\"ur relevante Neutronentrennenergie ist
        \[
            S_n(\isotope{16}{O})
            =E_B(\isotope{16}{O})-E_B(\isotope{15}{O}).
        \]
        Einsetzen ergibt
        \[
            S_n(\isotope{16}{O})
            =\SI{127.62}{MeV}-\SI{111.96}{MeV}
            =\SI{15.66}{MeV}.
        \]
        Dieses Neutron stammt aus der h\"ochsten besetzten Neutronenschale, also aus $1p_{1/2}$.
    \end{calculationbox}

    \begin{calculationbox}
        Nun f\"ugen wir gedanklich ein Neutron zu $\isotope{16}{O}$ hinzu. Dann entsteht $\isotope{17}{O}$. Die Bindung dieses zus\"atzlichen Neutrons ist
        \[
            S_n(\isotope{17}{O})
            =E_B(\isotope{17}{O})-E_B(\isotope{16}{O}).
        \]
        Einsetzen ergibt
        \[
            S_n(\isotope{17}{O})
            =\SI{131.76}{MeV}-\SI{127.62}{MeV}
            =\SI{4.14}{MeV}.
        \]
        Dieses zus\"atzliche Neutron sitzt in der n\"achsten freien Neutronenschale $1d_{5/2}$.
    \end{calculationbox}

    \begin{calculationbox}[title=Schalenabstand]
        Die Einteilchenenergien sind n\"aherungsweise negativ zu den Trennenergien:
        \[
            \epsilon(1p_{1/2})\approx -S_n(\isotope{16}{O}),
            \qquad
            \epsilon(1d_{5/2})\approx -S_n(\isotope{17}{O}).
        \]
        Daher ist der Abstand
        \[
            \Delta E
            =\epsilon(1d_{5/2})-\epsilon(1p_{1/2})
            \approx S_n(\isotope{16}{O})-S_n(\isotope{17}{O}).
        \]
        Also
        \[
            \Delta E
            =\SI{15.66}{MeV}-\SI{4.14}{MeV}
            =\SI{11.52}{MeV}.
        \]
    \end{calculationbox}

    \begin{answerbox}
        \[
            \boxed{\Delta E(1p_{1/2}\to 1d_{5/2})\approx \SI{11.52}{MeV}}.
        \]
    \end{answerbox}

    \item \textbf{Interpretation von $E_B(\isotope{17}{O})-E_B(\isotope{17}{F})$.}

    \begin{calculationbox}
        Zuerst bilden wir die Differenz:
        \[
            E_B(\isotope{17}{O})-E_B(\isotope{17}{F})
            =\SI{131.76}{MeV}-\SI{128.22}{MeV}
            =\SI{3.54}{MeV}.
        \]
    \end{calculationbox}

    \begin{solutionbox}
        Die beiden Kerne sind Spiegelkerne:
        \[
            \isotope{17}{O}: Z=8,\;N=9,
            \qquad
            \isotope{17}{F}: Z=9,\;N=8.
        \]
        Sie entstehen aus dem doppeltmagischen Kern $\isotope{16}{O}$ durch Hinzuf\"ugen eines Nukleons:
        \[
            \isotope{16}{O}+n\to\isotope{17}{O},
            \qquad
            \isotope{16}{O}+p\to\isotope{17}{F}.
        \]
        Bei idealer Ladungssymmetrie der starken Wechselwirkung w\"aren die starken Kernkraftbeitr\"age sehr \"ahnlich. Der Hauptunterschied ist daher die Coulomb-Absto\ss{}ung: In $\isotope{17}{F}$ gibt es ein zus\"atzliches Proton, das von den anderen Protonen elektrisch abgesto\ss{}en wird. Dadurch ist $\isotope{17}{F}$ weniger stark gebunden.
    \end{solutionbox}

    \begin{answerbox}
        \[
            \boxed{E_B(\isotope{17}{O})-E_B(\isotope{17}{F})=\SI{3.54}{MeV}}.
        \]
        Der Unterschied wird haupts\"achlich als zus\"atzliche Coulomb-Energie des protonenreicheren Spiegelkerns $\isotope{17}{F}$ interpretiert.
    \end{answerbox}

    \item \textbf{Radiusabsch\"atzung aus der Coulomb-Energie.}

    \begin{calculationbox}
        F\"ur eine homogen geladene Kugel mit abgezogener Selbstwechselwirkung der Protonen gilt laut Angabe
        \[
            E_C(Z)=\frac{3}{5}\frac{Z(Z-1)\alpha\hbar c}{R}.
        \]
        F\"ur $\isotope{17}{F}$ ist $Z=9$, f\"ur $\isotope{17}{O}$ ist $Z=8$. Daher ist die Coulomb-Energiedifferenz
        \[
            \Delta E_C
            =E_C(Z=9)-E_C(Z=8).
        \]
        Einsetzen liefert
        \[
            \Delta E_C
            =\frac{3}{5}\frac{\alpha\hbar c}{R}
            \bracket{9\cdot 8-8\cdot 7}.
        \]
        Der Klammerausdruck ist
        \[
            9\cdot 8-8\cdot 7=72-56=16.
        \]
        Also
        \[
            \Delta E_C
            =\frac{3}{5}\frac{16\alpha\hbar c}{R}
            =\frac{48}{5}\frac{\alpha\hbar c}{R}.
        \]
    \end{calculationbox}

    \begin{calculationbox}[title=Einsetzen der Bindungsenergiedifferenz]
        Wir identifizieren n\"aherungsweise
        \[
            \Delta E_C \approx E_B(\isotope{17}{O})-E_B(\isotope{17}{F})=\SI{3.54}{MeV}.
        \]
        Damit
        \[
            \SI{3.54}{MeV}
            =\frac{48}{5}\frac{\alpha\hbar c}{R}.
        \]
        Aufl\"osen nach $R$ ergibt
        \[
            R=\frac{48}{5}\frac{\alpha\hbar c}{\SI{3.54}{MeV}}.
        \]
        Mit $\alpha\hbar c\approx\SI{1.440}{MeV.fm}$ folgt
        \[
            R
            =\frac{48}{5}\frac{\SI{1.440}{MeV.fm}}{\SI{3.54}{MeV}}
            \approx \SI{3.91}{fm}.
        \]
    \end{calculationbox}

    \begin{answerbox}
        \[
            \boxed{R_{\mathrm{Coulomb}}\approx\SI{3.9}{fm}}.
        \]
    \end{answerbox}

    \item \textbf{Vergleich mit $R=R_0A^{1/3}$.}

    \begin{calculationbox}
        Mit $A=17$ und $R_0=\SI{1.2}{fm}$ gilt
        \[
            R_{0}A^{1/3}
            =\SI{1.2}{fm}\cdot 17^{1/3}.
        \]
        Numerisch ist
        \[
            17^{1/3}\approx 2.571.
        \]
        Also
        \[
            R_{0}A^{1/3}
            \approx \SI{1.2}{fm}\cdot 2.571
            \approx \SI{3.09}{fm}.
        \]
        Der relative Unterschied ist
        \[
            \frac{R_{\mathrm{Coulomb}}-R_0A^{1/3}}{R_0A^{1/3}}
            \approx
            \frac{\SI{3.91}{fm}-\SI{3.09}{fm}}{\SI{3.09}{fm}}
            \approx 0.27.
        \]
    \end{calculationbox}

    \begin{solutionbox}
        Die Coulomb-Absch\"atzung liefert einen gr\"o\ss{}eren Radius:
        \[
            R_{\mathrm{Coulomb}}\approx\SI{3.9}{fm}
            \qquad\text{gegen\"uber}\qquad
            R_0A^{1/3}\approx\SI{3.1}{fm}.
        \]
        Die Abweichung ist nicht \"uberraschend, weil die Rechnung stark idealisiert ist:
        \begin{itemize}
            \item Der Kern wurde als homogen geladene Kugel modelliert. Reale Kerne haben eine diffuse Randzone.
            \item Die Bindungsenergiedifferenz wurde vollst\"andig der Coulomb-Energie zugeschrieben. Kleine Beitr\"age durch Schalenstruktur, Isospinbrechung und unterschiedliche Einteilchenwellenfunktionen werden vernachl\"assigt.
            \item $R_0A^{1/3}$ ist eine globale Faustformel. Sie ist gut f\"ur Gr\"o\ss{}enordnungen, aber nicht f\"ur genaue Spiegelkern-Differenzen.
        \end{itemize}
    \end{solutionbox}

    \begin{answerbox}
        \[
            \boxed{R_{\mathrm{Coulomb}}\approx\SI{3.9}{fm}},
            \qquad
            \boxed{R_0A^{1/3}\approx\SI{3.1}{fm}}.
        \]
        Die Coulomb-Absch\"atzung ist etwa $27\,\%$ gr\"o\ss{}er. Die Gr\"o\ss{}enordnung stimmt, aber das Modell ist zu einfach f\"ur eine pr\"azise Radiusbestimmung.
    \end{answerbox}
\end{enumerate}

\subsection{Tafelvortrag zu Aufgabe 2}

\begin{talkbox}
\textbf{Tafelstruktur in 5 Schritten:}
\begin{enumerate}
    \item $\isotope{16}{O}$ ist doppeltmagisch: $Z=N=8$. Die Schale $1p_{1/2}$ ist voll.
    \item Neutron entfernen:
    \[
        S_n(\isotope{16}{O})=127{,}62-111{,}96=\SI{15.66}{MeV}.
    \]
    \item Neutron hinzuf\"ugen:
    \[
        S_n(\isotope{17}{O})=131{,}76-127{,}62=\SI{4.14}{MeV}.
    \]
    \item Schalenabstand:
    \[
        \Delta E=S_n(\isotope{16}{O})-S_n(\isotope{17}{O})=\SI{11.52}{MeV}.
    \]
    \item Spiegelkerne:
    \[
        E_B(\isotope{17}{O})-E_B(\isotope{17}{F})=\SI{3.54}{MeV}.
    \]
    Diese Differenz kommt haupts\"achlich von der Coulomb-Absto\ss{}ung im protonenreicheren $\isotope{17}{F}$.
    \[
        \SI{3.54}{MeV}=\frac{3}{5}\frac{\alpha\hbar c}{R}\bracket{9\cdot 8-8\cdot 7}
    \]
    und damit
    \[
        R\approx\SI{3.9}{fm}.
    \]
    Vergleich:
    \[
        R_0A^{1/3}=\SI{1.2}{fm}\cdot 17^{1/3}\approx\SI{3.1}{fm}.
    \]
\end{enumerate}
\end{talkbox}

\FloatBarrier
\newpage

%%%%%%%%%%%%%%%%%%%%%%%%%%%%%%%%%%%%%%%%%%%%%%%%%%%%%%%%%%%%%%%%%%%
\section{Gesamtzusammenfassung f\"ur die Pr\"ufung}

\begin{summarybox}[title={Formeln, die man k\"onnen sollte}]
\begin{align*}
    \vect{j}&=\vect{l}+\vect{s},\\
    \vect{l}\cdot\vect{s}&=\half\paren{\vect{j}^{\,2}-\vect{l}^{\,2}-\vect{s}^{\,2}},\\
    \avg{\vect{l}\cdot\vect{s}}&=\half\bracket{j(j+1)-l(l+1)-s(s+1)},\\
    S_n(A,Z)&=E_B(A,Z)-E_B(A-1,Z),\\
    E_C(Z)&=\frac{3}{5}\frac{Z(Z-1)\alpha\hbar c}{R},\\
    R&\approx R_0A^{1/3}.
\end{align*}
\end{summarybox}

\begin{summarybox}[title=Endergebnisse]
\begin{itemize}
    \item $1d$ bedeutet $l=2$.
    \item Aufspaltung: $1d_{5/2}$ und $1d_{3/2}$.
    \item Erwartungswerte:
    \[
        \avg{\vect{l}\cdot\vect{s}}_{1d_{5/2}}=1,
        \qquad
        \avg{\vect{l}\cdot\vect{s}}_{1d_{3/2}}=-\frac{3}{2}.
    \]
    \item Energiedifferenz:
    \[
        E_{5/2}-E_{3/2}=\frac{5}{2}V_{ls}(r).
    \]
    \item Schalenabstand in der Umgebung von $\isotope{16}{O}$:
    \[
        \Delta E(1p_{1/2}\to 1d_{5/2})\approx\SI{11.52}{MeV}.
    \]
    \item Spiegelkern-Bindungsenergiedifferenz:
    \[
        E_B(\isotope{17}{O})-E_B(\isotope{17}{F})=\SI{3.54}{MeV}.
    \]
    \item Radius aus Coulomb-Absch\"atzung:
    \[
        R\approx\SI{3.9}{fm}.
    \]
    \item Radius aus Faustformel:
    \[
        R\approx\SI{3.1}{fm}.
    \]
\end{itemize}
\end{summarybox}

\begin{warningbox}[title=Typische Fehler]
\begin{itemize}
    \item In Aufgabe 1 darf man nicht $\vect{j}=\vect{k}+\vect{s}$ schreiben; korrekt ist $\vect{j}=\vect{l}+\vect{s}$.
    \item Beim Quadrieren gilt $(\vect{l}+\vect{s})^2=\vect{l}^{\,2}+2\vect{l}\cdot\vect{s}+\vect{s}^{\,2}$, nicht $(\vect{l}^{\,2}+\vect{s}^{\,2})^2$.
    \item F\"ur $j=3/2$ ergibt sich $\avg{\vect{l}\cdot\vect{s}}=-3/2$, nicht $-1$.
    \item In Aufgabe 2(a) muss $E_B(\isotope{16}{O})=\SI{127.62}{MeV}$ verwendet werden.
    \item Bei der Coulomb-Absch\"atzung ist der Faktor $Z(Z-1)$ wichtig, weil die Selbstwechselwirkung der Protonen abgezogen wurde.
\end{itemize}
\end{warningbox}

\end{document}
